% ======================================================================
\section{Limits, consistency checks and comparison with the literature}\label{sec:checks}
% ======================================================================
% Sources: paper/notes/README.md, notes/lit_dglap_rc.tex, notes/lit_finite_dilute.tex,
% inputs/reviews/review_2.txt (secs. 3-6, 11, 13), limits_checked/*.
% Items marked with \pending{...} depend on the open fixes P1-P6 (BRIEF_fix.md) and must be
% revised when sections dglap_rc and finite are filled in.

The calculation involves a large number of terms, with non-commuting charges and Wilson lines in a definite order. In this section we list the checks to which the result has been subjected and compare it with known results. Unless stated otherwise, the algebraic identities quoted below were verified both analytically and numerically: the expressions were evaluated on a transverse lattice with random $SU(2)$ and $SU(3)$ Wilson lines in the adjoint representation and with either a completely general ordered charge correlator $\langle\rho^b(\v u)\rho^c(\v v)\rangle$ or the colour-singlet one, Eq.~\eqref{eq:singlet}. Since the Wilson lines are random, an identity that holds on the lattice to round-off error cannot hold by accident.

\subsection{Limits}\label{sec:limits}

\paragraph{No target, pure-gauge target.} For $U=1$ the operator $\bar\Omega^\dagger\Omega$ is the identity and every amplitude vanishes. In the result this is manifest: every term of every row, and every term of the BFKL, JIMWLK, DGLAP and running-coupling contributions, contains a factor that vanishes at $U=1$, such as $[U(\v x')-U(\v w')]^{ab'}\rho^{b'}(\v x')$, its analogue with three Wilson lines, or a combination of $f^{abc}$ and Wilson lines that reduces to $N_c-N_c-N_c+N_c=0$. A constant $U$ is a global colour rotation, which cannot produce gluons; all terms vanish for constant $U$ as well, by $f^{abc}U^{bd}U^{ce}=U^{ah}f^{hde}$ and $U^TU=1$.

\paragraph{Counting of charges.} The BFKL and JIMWLK terms are of order $g^4\rho^2$. No term of order $g^4\rho^4$ survives (Section~\ref{sec:fourrho}), and the three-$\rho$ terms that multiply a rapidity logarithm cancel in the symmetrised sense (Section~\ref{sec:threerho}). The finite terms are of order $g^4\rho^2$ and $g^4\rho^3$.

\paragraph{Reality.} Groups I and III are written with their complex conjugates and are real. Group II is Hermitian by itself: its kernels are symmetric under $(\v x,\v w)\leftrightarrow(\v x',\v w')$ and so is its colour operator, which we verified numerically to $10^{-12}$.

\paragraph{Odd terms.} Terms of the form $if^{abc}\times$(real kernel)$\times$(real colour structure)$\times e^{-i\v k\cdot(\v w-\v w')}$ give, together with their complex conjugate, $2\sin[\v k\cdot(\v w-\v w')]\times$(real), which is odd under $\v k\to-\v k$. All three-$\rho$ terms are of this type; they vanish for a Gaussian (colour-singlet, McLerran--Venugopalan-type) average over the projectile, which has no odd moments. This is the behaviour of the ``first saturation correction'' found in~\cite{McLerran:2016snu,Kovchegov:2018jun}: odd in $\v k$ configuration by configuration and zero after the average for the single inclusive spectrum. \pending{To be restated with the colour-averaged finite terms of Section~\ref{sec:finite} (input P6).}

\paragraph{Scaling.} Under a dilatation of all transverse coordinates, $\v x\to\lambda\v x$, $\v k\to\v k/\lambda$, every finite term scales like the LO spectrum. The finite terms contain no renormalisation scale; the dependence on $\mu$ is confined to the collinear and running-coupling terms (Section~\ref{sec:dglap_rc}).

\paragraph{Dilute target.} For $U=1+O(\alpha)$ the LO spectrum reduces to the $k_T$-factorised form, Eq.~\eqref{eq:kTfact}, with the Lipatov-type vertex. In this limit the BFKL and JIMWLK terms reduce to the LO BFKL evolution of the projectile and of the target unintegrated gluon distributions, respectively. A term-by-term comparison of the finite terms with the one-loop gluon production vertex in $k_T$ factorisation~\cite{Fadin:1998py,DelDuca:1998cx,Bartels:2006hg} requires all finite terms in momentum space and has not been carried out; see Section~\ref{sec:lit_constant} for the comparison of the constant.

\subsection{Leading order}

The LO spectrum, Eq.~\eqref{eq:LOsinglet}, coincides with the result of Kovchegov and Tuchin~\cite{Kovchegov:2001sc}, including its normalisation: for a quark projectile, $n\to C_F$, the prefactor is $\alpha_sC_F/(4\pi^4)$. The target operator $\mathcal O_{\rm LO}$ is the one of Kovner, Lublinsky and Weigert~\cite{Kovner:2006wr,Kovner:2006ge}.

\subsection{Projectile evolution: the BFKL term}\label{sec:check_BFKL}

\paragraph{Internal consistency.} The ordered form~\eqref{eq:BFKLrho} is an exact regrouping of the term-by-term list of all two-$\rho$ terms proportional to $\ln(\vee/k^+)$, for an arbitrary ordered correlator (relative difference $\lesssim10^{-15}$ for $SU(2)$, $10^{-14}$ for $SU(3)$), and Eq.~\eqref{eq:BFKLsinglet} is its colour-singlet reduction to the same accuracy.

\paragraph{Exact one-gluon evolution of a colour-singlet projectile.} As an independent test we computed the change of the spectrum due to one additional gluon with $k^+<p^+<\vee$ in a colour-singlet projectile directly, without the expansion in $\rho$ used in the rest of the paper. Any set of colour sources $s$ at positions $\v x_s$, with charges $Q_s^e$, produces the observed gluon with the operator $\sum_{s,s'}Q^e_{s'}Q^e_sF(\v x_s,\v x_{s'})$, where
\begin{align}
F(\v a,\v b)=\int_{w,w'}e^{-i\v k\cdot(\v w-\v w')}\Kv(\v a-\v w)\cdot\Kv(\v b-\v w')\,{\rm Tr}\big[(U^\dagger(\v a)-U^\dagger(\v w))(U(\v b)-U(\v w'))\big]\,,
\end{align}
since the colour rotation of the sources cancels in the inclusive sum. For a quark--antiquark dipole at $\v x_1$, $\v x_2$ the three charges of the $q\bar qg$ colour singlet satisfy $Q_1^2=Q_2^2=C_F$, $Q_g^2=N_c$, $Q_1\cdot Q_2=1/(2N_c)$ and $Q_{1,2}\cdot Q_g=-N_c/2$, and the probability to emit the additional gluon at $\v z$ is $\frac{g^2}{4\pi^3}\ln\frac{\vee}{k^+}C_F\,K_D(\v x_1,\v x_2,\v z)$. Real minus virtual gives, exactly at finite $N_c$,
\begin{align}
\delta\Big[\frac{d^3N}{d^3k}\Big]&=\frac{\alpha_sN_c}{2\pi^2}\ln\frac{\vee}{k^+}\int_zK_D(\v x_1,\v x_2,\v z)\big[\mathcal D_{1z}+\mathcal D_{z2}-\mathcal D_{12}\big]\times(\text{LO prefactor})\,,\notag\\
\mathcal D_{ab}&=\frac{1}{2N_c}\big[F_{aa}+F_{bb}-F_{ab}-F_{ba}\big]\,,
\end{align}
the dipole BFKL (linearised BK) kernel acting on the projectile. Substituting $n=C_F[\delta_{11}+\delta_{22}-\delta_{12}-\delta_{21}]$ into Eq.~\eqref{eq:BFKLsinglet} reproduces this result pointwise in $\v z$. Since every symmetric $n(\v x,\v y)$ with vanishing row sums, i.e.\ every colour-singlet projectile, is a linear combination of such dipole correlators, and Eq.~\eqref{eq:BFKLsinglet} is linear in $n$, the BFKL term is the exact leading-logarithmic projectile evolution for every colour-singlet projectile. The calculation with fundamental charges is exact to all orders in the charges, which shows that the expansion in $\rho$ loses nothing for a dipole.

\paragraph{Literature.} This is the structure found by Kovner and Lublinsky for the single inclusive spectrum~\cite{Kovner:2006wr} and by Kovner, Lublinsky and Weigert~\cite{Kovner:2006ge}: the projectile correlator evolves with BFKL from the projectile rapidity down to the rapidity of the produced gluon. It is the first-order expansion of the leading-logarithmic factorisation theorem of Gelis, Lappi and Venugopalan~\cite{Gelis:2008rw}, and it satisfies the general structure of the NLO correction discussed by Gelis and Laidet~\cite{Gelis:2012ri}.

\subsection{Target evolution: the JIMWLK term}\label{sec:check_JIMWLK}

We computed $-H_{\rm JIMWLK}\,\mathcal O_{\rm LO}$ independently, with the Hamiltonian~\eqref{eq:HJIMWLK} in the normalisation of~\cite{Lublinsky:2016meo,Kovner:2013ona}, for a general ordered correlator. The term-by-term list of the two-$\rho$ terms proportional to $\ln(k^+/\Lambda)$, the grouped form~\eqref{eq:JIMWLKrho} and $-H_{\rm JIMWLK}$ acting on the LO spectrum agree to $10^{-14}$ for $SU(2)$ and $SU(3)$, and the colour-singlet form~\eqref{eq:JIMWLKsinglet} agrees with the colour-singlet projection to $10^{-15}$. The JIMWLK term is therefore the JIMWLK evolution of the LO spectrum, with the correct sign and normalisation, for an arbitrary projectile and not only for a colour-singlet one. Together with Section~\ref{sec:check_BFKL}, this confirms at fixed order the factorisation of the rapidity logarithms at the rapidity of the produced gluon~\cite{Kovner:2006wr,Kovner:2006ge,Gelis:2008rw}.

\subsection{Wave function}\label{sec:check_wf}

The light-cone wave function from which the coefficients of $\Omega$ are obtained was computed to order $g^3$, with the same longitudinal cut-offs ($\vee=\Lambda e^{\delta Y}$) and $d_\perp=2-\epsilon$, by Lublinsky and Mulian~\cite{Lublinsky:2016meo}. Our components agree with theirs as follows:
\begin{center}
\renewcommand{\arraystretch}{1.2}
\small
\begin{tabular}{@{}p{4.6cm}p{2.3cm}p{6.6cm}@{}}
\toprule
component & Ref.~\cite{Lublinsky:2016meo} & comparison\\
\midrule
LO soft gluon, normalisation $\mf N_2$ & eqs.~2.50--2.52 & identical\\
two gluons, $O(g^2)$ & eqs.~3.95--3.96 & identical\\
one gluon, $O(g^3)$, linear in $\rho$ & eqs.~3.97--3.98 & identical for $N_f=0$, incl.\ $\frac{11}{3}$, $\frac{67}{9}-\frac{\pi^2}{3}$, $\ln^2$ cut-off terms\\
one gluon, $O(g^3)$, $f^{abc}\{\rho^b,\rho^c\}$ & eqs.~3.100--3.101 & equal after the antisymmetrisation enforced by $f^{abc}\{\rho^b,\rho^c\}$\\
one gluon, $O(g^3)$, $\rho\rho\rho$ & eq.~3.102 & identical up to the operator ordering, see below\\
\bottomrule
\end{tabular}
\end{center}
In the two-$\rho$ component the two forms differ by terms symmetric under the exchange of the two charges, which do not contribute. The coordinate-space forms of the $\{\rho,\rho\}$ component, Eq.~\eqref{eq:A3tworho}, and of the coefficients go beyond the published material. The operator ordering of the three-$\rho$ component, $\rho^a(\v w)\rho^b(\v x)\rho^b(\v y)$, i.e.\ $[\mf A^{(3)}]_{\rho\rho\rho}=\mf A\,\mf N_2$, follows from third-order perturbation theory with operator-valued matrix elements; we verified it against a construction of $\Omega$ and an exact diagonalisation with finite-dimensional matrices representing the charges (agreement to $3\times10^{-16}$; the opposite ordering deviates at order one).

\subsection{The constant \texorpdfstring{$\frac{67}{9}-\frac{\pi^2}{3}$}{67/9-pi2/3}}\label{sec:lit_constant}

The term of Row I.II proportional to the LO structure is, with its complex conjugate,
\begin{align}\label{eq:constantLO}
\frac{\alpha_sN_c}{2\pi}\Big(\frac{67}{9}-\frac{\pi^2}{3}\Big)\frac{d^3N_{\rm LO}}{d^3k}\,.
\end{align}
It follows directly from the one-$\rho$ part of $\mf A^{(3)}$, Eq.~\eqref{eq:A3onerho}, and the normalisation of the LO amplitude. Before the cancellation of the ultraviolet $\beta_0$ logarithms against Row I.V, Row I.II with its complex conjugate is $\frac{\alpha_sN_c}{2\pi}\big[\frac{11}{3}\big(2\gamma_E+\frac12\ln\frac{X^2\mu^2_{\MS}}{4}+\frac12\ln\frac{X'^2\mu^2_{\MS}}{4}\big)+\frac{67}{9}-\frac{\pi^2}{3}\big]\times$LO, $X=x-w$, $X'=x'-w'$. This is, term by term at $N_f=0$, the bare one-gluon kernel $K_{JSJ}$ of the NLO JIMWLK Hamiltonian in the form of Kovner, Lublinsky, Skokov and Zhao~\cite{Kovner:2023vsy}, with the same sign, size and scale; since both are built from the wave function of~\cite{Lublinsky:2016meo}, this is a consistency check rather than an independent comparison. The inclusive quantities in which the same constant appears, the NLO BK kernel of Balitsky and Chirilli~\cite{Balitsky:2007feb}, the NLO BFKL kernel of Fadin and Lipatov~\cite{Fadin:1998py} and the NLO jet vertex of Bartels, Sabio Vera and Schwennsen~\cite{Bartels:2006hg}, have the coefficient $\frac{\alpha_sN_c}{4\pi}(\frac{67}{9}-\frac{\pi^2}{3})$, half of Eq.~\eqref{eq:constantLO}. These quantities are inclusive over the collinear gluon pair; as noted in~\cite{Kovner:2023vsy}, half of the constant sits there in the two-gluon terms. For the single inclusive gluon, whose collinear splitting is factorised into the fragmentation function (Section~\ref{sec:dglap_rc}), the bare-vertex value~\eqref{eq:constantLO} is the appropriate one, and the other half must reside in the collinear finite terms. \pending{Not verified: the momentum-weighted ($\int d\xi\,\xi$) inclusive combination of all collinear terms should give $-\frac{\alpha_sN_c}{4\pi}(\frac{67}{9}-\frac{\pi^2}{3})\times$LO; depends on the final form of the collinear and finite terms (P3, P4).} The amplitude-level one-loop Lipatov vertex~\cite{DelDuca:1998cx} contains no $\frac{67}{9}$ (it appears only after the real two-gluon production is added) and cannot be compared term by term with the light-cone wave-function constant.

\subsection{Collinear and running-coupling terms}\label{sec:lit_dglap}
% Depends on P3, P4 (input notes/fix_poles_dglap.tex). Revise after sec_dglap_rc is filled.

\pending{This subsection summarises the comparison of notes/lit\_dglap\_rc.tex and must be revised when Section~\ref{sec:dglap_rc} is written (P3, P4).} The collinear term comes from Row II.IV, in which the observed gluon is produced in the splitting $g\to gg$ of a gluon of momentum $k^+/\xi$ that was produced at LO; it is a final-state splitting. Compared with the forward hadron production results of Chirilli, Xiao and Yuan~\cite{Chirilli:2011km,Chirilli:2012jd} and of Altinoluk et al.~\cite{Altinoluk:2014eka}:
\begin{itemize}
\item The splitting function has the normalisation of the LO $P_{gg}$, $2N_c[\frac{\xi}{(1-\xi)_+}+\frac{1-\xi}{\xi}+\xi(1-\xi)]+\frac{\beta_0}{2}\delta(1-\xi)$ at $N_f=0$, and the coefficient $\beta_0/2$ of $\delta(1-\xi)$ is the one required by the momentum sum rule $\int_0^1d\xi\,\xi P_{gg}(\xi)=0$. The prefactor $\frac{\alpha_s}{2\pi}$, the factor $\xi^{-2}$ with the phase $e^{-i\v k\cdot(\v y-\v y')/\xi}$, the sign and the dependence on $\mu$ are those of a final-state $\MS$ fragmentation of the LO spectrum of the parent gluon.
\item The end point $\xi\to1$ (integrated gluon soft) is plus-subtracted into the JIMWLK term, as the rapidity divergence is subtracted into BK evolution in~\cite{Chirilli:2012jd}. The end point $\xi\to0$ (observed gluon soft with respect to its parent) is subtracted into the BFKL term; the combination of the two subtractions is exactly the ``double plus'' prescription of~\cite{Kovner:2023vsy}, for the same reason: both poles are part of the LL evolution. In the forward (hybrid) case $\xi$ is bounded from below and no subtraction at $\xi\to0$ arises.
\item The argument of the collinear logarithm is the transverse separation of the observed gluon in the amplitude and in the conjugate amplitude, $(\v y-\v y')/\xi$, whereas in the $\MS$ form of~\cite{Chirilli:2012jd,Altinoluk:2014eka} it is the separation of the parent, $\v y-\v y'$. The difference is a finite term $\int d\xi\,P_{gg}(\xi)\ln\xi^2\otimes$LO, which depends on the $d$-dimensional definition of the convolution. Our form coincides with the cut-off scheme of~\cite{Altinoluk:2014eka}, with $\mu^2_{\MS}=c_0^2\mu^2_{\rm cut}\approx1.26\,\mu^2_{\rm cut}$, $c_0=2e^{-\gamma_E}$. \pending{P4: final form of the logarithm and of the $\xi\to0$ subtraction.}
\item No counterpart arises to the initial-state DGLAP evolution of the projectile parton distribution of the hybrid formalism: the projectile here is described by an unintegrated charge correlator, whose small-$x$ evolution is the BFKL term.
\item The $\beta_0$ terms. The real emission produces no $\delta(1-\xi)$ term, and the virtual $\frac{11}{3}$ terms of the vertex correction (Row I.II) cancel against those of Row I.V. The term $\frac{\beta_0}{2}\delta(1-\xi)$ of $P_{gg}$ is added and subtracted, the subtracted term being the running-coupling term; with it the LO coupling becomes $g^2(\mu)$ at the scale of the collinear subtraction, or $g^2(c_0^2/r^2)$ with $r$ the transverse separation conjugate to $\v k$. In the classification of~\cite{Kovner:2023vsy}, where of the $2\beta_0$ of the NLO JIMWLK kernel one is the ultraviolet running of the coupling and one a DGLAP-type resolution logarithm, our $\beta_0$ terms are of the second, collinear, type. This agrees with Kovchegov and Weigert~\cite{Kovchegov:2007vf} and with Horowitz and Kovchegov~\cite{Horowitz:2010yg}, who find that the coupling of the emission vertex of the produced gluon runs with the collinear resolution scale, $\alpha_s(\Lambda^2_{\rm coll})$, which plays the role of the factorisation scale of the fragmentation function. The other couplings of~\cite{Horowitz:2010yg}, associated with the $t$-channel exchanges, are not renormalised in this calculation. \pending{P3: account of all $1/\epsilon$ poles; wording of the running-coupling term.}
\end{itemize}

\subsection{Finite terms}\label{sec:lit_finite}
% Depends on P1, P2, P5, P6. Revise after sec_finite is filled.

\pending{To be completed after Section~\ref{sec:finite} (P1, P2, P5, P6).} The following properties of the finite terms are established independently of the open points: (i) all finite terms vanish for $U=1$ and for constant $U$; (ii) Group II is real without its complex conjugate; (iii) in Group II, Row I, individual terms have denominators that vanish on whole surfaces in the transverse plane (e.g.\ $(\v x'-\v w')^2(\v x-\v z)^2=(\v x-\v w)^2(\v x'-\v z)^2$), but the sum is finite and continuous across them, so that no principal-value prescription is needed there; this is the mechanism known from the kernels $\ln(X^2Y'^2/X'^2Y^2)/(X^2Y'^2-X'^2Y^2)$ of NLO JIMWLK~\cite{Balitsky:2013fea,Kovner:2013ona}; (iv) the kernels are integrable in all coincidence limits of Groups II and III and of Group I Row V; (v) the collinear limit of Row II.IV agrees with collinear factorisation (Section~\ref{sec:lit_dglap}), and its limits $\xi\to0$ and $\xi\to1$ reproduce the coefficients of the projectile-side and target-side soft-gluon emission, as in the corresponding limits of the NLO jet vertex~\cite{Bartels:2006hg}. The comparison of the full dilute limit with the jet vertex of~\cite{Bartels:2006hg,Fadin:1998py} is structural only. Open at the time of writing are the large-distance behaviour of the soft remainders in Group I Row III and Group III Rows III, V, VI, and the large-distance tail of Row II.IV for a colour-neutral projectile, $\propto\beta_0\frac{4\pi}{\v k^2}\int\frac{d^2Y}{Y^2}[-\frac{11}{6}+\frac16\cos2\varphi]$, for which no compensating term has been identified. The kinematic-constraint remainders $[k^+/(p^+(p^++k^+))]_+$ of Group III are the analogue of the terms discussed in the forward NLO literature~\cite{Stasto:2013cha,Altinoluk:2014eka,Ducloue:2016shw}.
